\section{Comparative Tests and Failure Conditions}
\subsection{Three-arm target-assignment experiment}
Randomize clinicians, researchers, or trained participants to identical longitudinal cases under three vocabularies: umbrella language (Arm U), experience-centered decomposition (Arm E), and SCTC decomposition (Arm $\Gamma$). All arms must answer the same forced response fields:
\begin{enumerate}
\item individual status: same individual / ended / indeterminate;
\item phenomenal experience: present / absent / unknown;
\item responsiveness: present / absent / unknown;
\item connectedness: present / absent / unknown.
\end{enumerate}
The missingness mechanism must be preregistered. ``Unknown'' is not scored as a reclassification.

Let $V_R$ be the number of unsupported flips in the system-level classification caused solely by changes in evidence availability, and let $\mathrm{Err}_E(R)$ be error on separately scored phenomenal-state judgments. The distinctive SCTC prediction is conjunctive:
\begin{equation}
\boxed{V_\Gamma<V_E\quad\land\quad \mathrm{Err}_E(\Gamma)-\mathrm{Err}_E(E)\le\varepsilon_E.}
\end{equation}
Reduced volatility is not a victory if it is purchased by making people worse at reasoning about experience.

\subsection{Closed-inventory system-level discontinuity test}
To prevent post hoc absorption of counterexamples, the faculty inventory is frozen in advance as $\mathbf F_B^*=(W,E,K,R,M,P,Q,S,I)$. Search across at least two preregistered paradigms for a variable $Z$ such that:
\begin{enumerate}
\item $Z$ changes while $L_B=1$;
\item $Z$ is not reducible to $\mathbf F_B^*$ or measurement availability;
\item $Z$ adds independent predictive or interventional information after those variables are controlled;
\item the relevant outcome family and reduction criteria were fixed before observing $Z$.
\end{enumerate}
If such a system-level $Z$ is necessary to recover the classifications science intends by consciousness, then $C\equiv L_B$ is insufficient and SCTC fails at its core.

\subsection{Explicit failure conditions}
SCTC should be abandoned or revised if: (i) the direct assignment experiment favors $\mathcal R_E$ or another rival; (ii) a preregistered independent system-level $Z$ survives the discontinuity test; (iii) biological individuality cannot be operationalized with adequate reliability for the intended domain; or (iv) the biological substrate restriction fails because a nonbiological system satisfies the deeper organizational criterion that the theory actually requires.

\section{What SCTC Does Not Entail}
Because the paper uses \emph{consciousness} at the minimal system level, several non-entailments must be explicit:
\begin{equation}
C_B=1\;\not\!\Rightarrow\;\substack{E=1,\;\text{pain, sentience, interests, personhood,}\\
\text{moral status, or legal status}.}
\end{equation}
Those are separate empirical or normative variables. The point is not to weaken them. It is to prevent them from being re-imported into $C$ after the theory has separated them.

The same rule applies to clinical language. ``Loss of consciousness'' should be replaced by the variable actually measured: loss of responsiveness, wakefulness, connectedness, memory formation, reportability, or evidence of phenomenal experience. ``Neural correlate of consciousness'' should identify the particular faculty or state whose correlate is being measured.

\section{Discussion: The Problem Was Allowed to Move}
The strongest criticism of SCTC is that experience was the target researchers cared about all along. That may often be true. It does not establish that experience and system-level consciousness are the same variable. It establishes that experience deserves its own theory.

This is the central methodological claim. A question can remain profound after a terminology correction. What changes is which variables inherit the mystery. If the unresolved question is why physical organization gives rise to subjective experience, then $E$ is the target. Study $E$ as $E$. Do not require wakefulness, memory, reportability, responsiveness, sensory access, or organismal continuity to inherit the same explanatory debt because one historical noun covered them all.

The paper therefore does not claim to have explained dreaming, pain, visual phenomenology, or the existence of a point of view. It claims that none of those unanswered questions can be used against Eq. \ref{eq:def} without first supplying an argument that the faculty in question is identical to the system-level variable.

\defbox{\textbf{The decisive question is not ``Have you explained experience?'' It is ``Why should an unexplained theory of experience count as an unexplained definition of the continuing system?''}}

That question cannot be answered by changing the meaning of consciousness halfway through the argument.

\section{Conclusion}
A scientific mechanism requires a stable target. Consciousness research inherited a word that has been used for a biological individual, its wakefulness, its responses, its access to the environment, its memories, its self-model, its reports, and its phenomenal experiences. Those variables do not move together.

SCTC fixes the system-level target first:
\begin{equation}
\boxed{\begin{aligned}
C_B(t)&\equiv L_B(t)\\
&=\text{viable organism-level existence of }B.
\end{aligned}}
\end{equation}
Everything else must earn its own variable.

This does not solve the hard problem of phenomenal experience. It prevents that problem from being silently substituted for every other question called consciousness. The anesthesia dream is the simplest demonstration. The person continues whether a dream occurs, does not occur, or cannot later be recovered. Why a dream occurred is a real scientific question about $E$. It is not evidence that the continuing individual disappeared and returned.

The historical hard problem may therefore contain two different things: a genuinely difficult problem of phenomenal experience and a prior representation error about what object the word \emph{consciousness} was supposed to denote. The first deserves explanation. The second deserves correction.

\begin{center}
\small\bfseries A mechanism can be mysterious only after its target is stable.
\end{center}

